% Balanced version (both routes with equal weight). Scrutiny log at the end.
% =====================================================================
% Shared preamble for draft.tex, draft-communication.tex, draft-interventional.tex
% Orange [FILL: ...] marks facts only you can supply. See skeleton.tex
% section 0.5 for the qualifiers that must survive editing.
% =====================================================================
\documentclass[11pt,a4paper]{article}
\usepackage[margin=2.2cm]{geometry}
\usepackage[backend=biber,style=numeric,sorting=nyt,maxbibnames=6,giveninits=true]{biblatex}
\usepackage{amsmath,amssymb,amsthm}
% \usepackage{amsthm}
% \usepackage[margin=2cm]{geometry}
\usepackage{dsfont}
\usepackage{tikz}
\usetikzlibrary{cd}
\usetikzlibrary{babel,calc}
\usetikzlibrary{positioning,arrows.meta,fit,backgrounds}
% \usepackage{tikz-cd}
% \usepackage{titlesec}
\usepackage{enumitem}
\usepackage{hyperref}
\usepackage{cleveref}
% \usepackage{parskip}
\usepackage{ebproof}
% \usepackage{quiver}
\usepackage{stackengine,scalerel}
% \usepackage{xparse}
% \usepackage{physics}
% \usepackage{physics2}
% \usephysicsmodule{ab,ab.braket}
% \usepackage{tikzit}
% \input{smc.tikzstyles}
\usepackage{xcolor}
\usepackage{graphicx}
\usepackage{array}    % m{} columns + \arraybackslash (combined car-model table)
\usepackage{booktabs} % \toprule/\midrule/\bottomrule

\usepackage{mathtools}
\usepackage{stmaryrd}
\usepackage{subcaption}

\newtheorem*{fact}{Fact}

\crefname{prop}{prop.}{props.}
\Crefname{prop}{Proposition}{Propositions} % Capitalized version for start of sentences

\crefname{remark}{remark}{remarks}
\Crefname{remark}{Remark}{Remarks} % Capitalized version for start of sentences

\crefname{appendix}{appendix}{appendices}
\Crefname{appendix}{Appendix}{Appendices} % Capitalized version for start of sentences

% \crefname{property}{property}{properties}
% \Crefname{property}{Property}{Properties} % Capitalized version for start of sentences

\crefformat{property}{#2#1#3}
\Crefformat{property}{#2#1#3}

\crefname{def}{definition}{definitions.}
\Crefname{def}{Definition}{Definitions} % Capitalized version for start of sentences

\crefformat{section}{#2\S#1#3}
\Crefformat{section}{#2\S#1#3}

\usepackage{listings}

\usepackage{tcolorbox}
\tcbuselibrary{listings} % Import listings support

\definecolor{codegreen}{rgb}{0,0.6,0}
\definecolor{codegray}{rgb}{0.5,0.5,0.5}
\definecolor{codepurple}{rgb}{0.58,0,0.82}
\definecolor{backcolour}{rgb}{0.95,0.95,0.92}

\lstdefinestyle{mystyle}{
    backgroundcolor=\color{backcolour},
    commentstyle=\color{codegreen}\textit,
    keywordstyle=\color{blue},
    numberstyle=\tiny\color{codegray},
    stringstyle=\color{codepurple},
    basicstyle=\ttfamily\footnotesize,
    breakatwhitespace=false,
    breaklines=false,
    captionpos=b,
    keepspaces=true,
    numbers=left,
    numbersep=5pt,
    showspaces=false,
    showstringspaces=false,
    showtabs=false,
    tabsize=2
}

\lstset{style=mystyle}

\tcbset{
    mystyle/.style={
        listing only,
        breakable, % Allows box to span pages
        colback=CodeBack,
        colframe=CodeFrame,
        boxrule=0.5pt,
        sharp corners,
        listing options={
            backgroundcolor=\color{backcolour},   
            commentstyle=\color{codegreen}\textit,
            keywordstyle=\color{blue},
            numberstyle=\tiny\color{codegray},
            stringstyle=\color{codepurple},
            basicstyle=\ttfamily\footnotesize,
            breakatwhitespace=false,         
            breaklines=true,                 
            captionpos=b,                    
            keepspaces=true,                 
            numbers=left,                    
            numbersep=5pt,                  
            showspaces=false,                
            showstringspaces=false,
            showtabs=false,                  
            tabsize=2,
        }
    }
}

\newtcblisting{codelisting}[3][]{
    listing only,
    % breakable, % Allows box to span pages
    colback=backcolour,
    % colframe=CodeFrame,
    boxrule=0pt,
    boxsep=0pt,
    top=1mm,
    bottom=1mm,
    sharp corners,
    notitle,
    listing options={
        language=#2,
        morekeywords= {exo, endo},
        commentstyle=\color{codegreen}\textit,
        keywordstyle=\color{blue},
        numberstyle=\tiny\color{codegray},
        stringstyle=\color{codepurple},
        basicstyle=\ttfamily\tiny,
        breakatwhitespace=false,         
        breaklines=true,                 
        captionpos=b,                    
        keepspaces=true,                 
        numbers=left,                    
        numbersep=5pt,                  
        showspaces=false,                
        showstringspaces=false,
        showtabs=false,                  
        tabsize=2,
        aboveskip=0pt,  % Removes the space above the code
        belowskip=0pt,  % Removes the space below the code
        mathescape=#3,
        #1               % optional extra listings options (e.g. escapeinside for \cnew)
    }
}

% Highlight code that is new relative to the previous rung's program.
% Use inside a codelisting that passes escapeinside, e.g.
%   \begin{codelisting}[escapeinside={(*}{*)}]{haskell}{true} ... (*\cnew{new}*) ...
\definecolor{codenew}{RGB}{255,228,130}
\newcommand{\cnew}[1]{{\setlength{\fboxsep}{0.6pt}\colorbox{codenew}{\ttfamily #1}}}
% Coloured (not highlighted) font, for marking new syntax in the grammar figure.
\newcommand{\cnewm}[1]{\textcolor{blue!65!black}{#1}}

% Theo's Macros
\newcommand{\naturals}{\mathbb{N}}
\newcommand{\integers}{\mathbb{Z}}
\newcommand{\booleans}{\mathbb{B}}
\newcommand{\indicator}{\mathds{1}}
\newcommand{\va}{\mathbf{a}}
\newcommand{\vb}{\mathbf{b}}
\newcommand{\vx}{\mathbf{x}}
\newcommand{\vy}{\mathbf{y}}
\newcommand{\vw}{\mathbf{w}}
\newcommand{\vz}{\mathbf{z}}
\newcommand{\vsigma}{\mathbf{\sigma}}
\newcommand{\reals}{\mathbb{R}}
\newcommand{\complex}{\mathbb{C}}
\newcommand{\field}{\mathbb{F}}
\newcommand{\sign}{\mathrm{sign}}
\DeclareMathOperator*{\argmin}{arg\,min}
\newcommand{\true}{\mathsf{True}}
\newcommand{\false}{\mathsf{False}}
\newcommand{\brackets}[1]{\{#1\}}

% \newcommand{\bicatb}{\mathcal{B}}

\newcommand{\catc}{\mathbf{C}}
\newcommand{\catcop}{\catc^{op}}
\newcommand{\catd}{\mathbf{D}}
\newcommand{\cate}{\mathcal{E}}
\newcommand{\eqdef}{\triangleq}
\newcommand{\id}{\mathbf{id}}
% \newcommand{\idbicat}{\mathbf{Id}}
\newcommand{\Id}{\mathbf{Id}}
\newcommand{\righttop}[1]{\xrightarrow{#1}}
\newcommand{\lefttop}[1]{\xleftarrow{#1}}
\newcommand{\Ob}{\mathbf{Ob} \ }
\newcommand{\Hom}{\mathrm{Hom}}
% \newcommand{\yo}{\textmc{よ}}
\newcommand{\yo}{\textbf{y}}

\newcommand{\terminalmorph}{\langle\rangle}
\newcommand{\prodmorph}[2]{\langle#1, #2\rangle}

\newcommand{\ul}[1]{\underline{#1}}
\newcommand{\ulyo}{\ul{\yo}}
\newcommand{\yozero}{\yo(0)}
\newcommand{\yoone}{\yo(1)}
\newcommand{\yon}{\yo(n)}
\newcommand{\yom}{\yo({m})}
\newcommand{\yop}{\yo(p)}
% \newcommand{\op}{\mathbf{op}}
\newcommand{\Ops}{\mathcal{O}}
% \newcommand{\typing}[3]{#1 \ | \ #2 \vdash #3}
% \newcommand{\typingfilled}[1]{\Gamma \ | \ \Delta \vdash #1}
% \newcommand\eye{\ensurestackMath{\stackinset{c}{}{c}{-.33pt}%
%   {\bullet}{\bigcirc}}}
% \newcommand\smallereye{\scalerel*{\eye}{x}}
% \newcommand{\eyemid}{ \ \eye \ }

% Vectors
\newcommand{\veca}{\Vec{a}}
\newcommand{\vecb}{\Vec{b}}
\newcommand{\vecc}{\Vec{c}}
\newcommand{\vecd}{\Vec{d}}
\newcommand{\vecx}{\Vec{x}}
\newcommand{\vecy}{\Vec{y}}
\newcommand{\vecz}{\Vec{z}}

\newcommand{\alphaequiv}{\simeq_\alpha}

\newcommand{\LHS}{\mathsf{LHS}}
\newcommand{\RHS}{\mathsf{RHS}}

% Chapter 2 algebraic theories
\newcommand{\Mon}{\mathsf{Mon}}

\newcommand{\Tree}{\mathsf{Tree}}

\newcommand{\Kl}[1]{\mathsf{Kl}({#1})}

\newcommand{\Name}{\mathcal{N}}
\newcommand{\Var}{\mathsf{Var}}
\newcommand{\IPVar}{\mathsf{IPVar}}
\newcommand{\GroundValue}{\mathsf{GroundValue}}
\newcommand{\Value}{\mathsf{Value}}
\newcommand{\Expr}{\mathsf{Expr}}

\newcommand{\bind}{\mathbin{>\!\!\!>\mkern-6.7mu=}}
\newcommand{\fin}{\mathsf{fin}}
\newcommand{\MVal}{\mathsf{MVal}}
\newcommand{\Intervention}{\mathsf{Intv}}
\newcommand{\MultiWorld}{\mathsf{MultiWorldCounterfactual}}
\newcommand{\TwinWorld}{\mathsf{TwinWorldCounterfactual}}
\newcommand{\supp}{\mathsf{supp}}

% ------------- Formal Def ----------------
\newcommand{\IdealisedChiRho}{\mathsf{ChiRho}}
\newcommand{\Prob}{\mathsf{Prob}}
\newcommand{\Dist}{\mathsf{Dist}}
\newcommand{\Paths}{\mathsf{Paths}}
\newcommand{\UPaths}{\mathsf{UPaths}}
\newcommand{\posreals}{\reals_{\geq 0}}

% Syntax
% \newcommand{\opif}[3]{\mathbf{if} \ #1 \ #2 \ #3}
% \newcommand{\lift}[1]{\overline{#1}}
\newcommand{\opprod}{\mathsf{prod}}
\newcommand{\opscore}{\mathsf{score}}
\newcommand{\opflip}{\mathsf{flip}}
\newcommand{\dbernoulli}{\mathbf{bernoulli}}
\newcommand{\dnormal}{\mathbf{normal}}
\newcommand{\constr}{\mathbf{constr}}
\newcommand{\intvnone}{\mathsf{None}}
\newcommand{\intvsome}{\mathsf{Some}}
\newcommand{\opskip}{\mathsf{skip}}
\newcommand{\opsample}{\mathsf{sample}}
\newcommand{\opintervene}{\mathsf{intervene}}
\newcommand{\opbroadcast}{\mathsf{broadcast}}
\newcommand{\optrack}{\mathsf{track}}
\newcommand{\opdistribute}{\mathsf{distribute}}
\newcommand{\opif}{\mathbf{if}}
\newcommand{\ifelse}[3]{\opif \ #1 \ \mathbf{then} \ #2 \ \mathbf{else} \ #3}
\newcommand{\opthunk}{\mathsf{thunk}}
\newcommand{\opforce}{\mathsf{force}}
\newcommand{\returnSymb}{\mathbf{return}}
\newcommand{\return}[1]{\returnSymb \ #1}
\newcommand{\returnParam}[2]{\returnSymb_{#1} \ #2}
\newcommand{\app}{\odot}
\newcommand{\pureval}{\mathbf{pureVal}}
\newcommand{\mem}{\mathsf{mem}}
\newcommand{\letbind}[3]{\mathbf{let} \ #1 = #2 \ \mathbf{in} \ #3}
\newcommand{\opdo}{\mathsf{do}}
\newcommand{\rhide}{\mathsf{hide}}
\newcommand{\rwk}{\mathsf{wk}}
\newcommand{\rcontr}{\mathsf{contr}}
\newcommand{\rexch}{\mathsf{exch}}
\newcommand{\rrename}{\mathsf{rename}}
\newcommand{\Renamings}{\mathsf{Renamings}}
\newcommand{\Deltahat}{\widehat{\Delta}}
\newcommand{\projf}{\pi_{\mathsf{f}}}
\newcommand{\projcf}{\pi_{\mathsf{cf}}}
\newcommand{\coerce}[3]{\mathsf{coerce}_{#2}^{#1}[#3]}
\newcommand{\rv}{\underline{r}}
\newcommand{\restr}{\rvert}
\newcommand{\subst}[2]{[#1/#2]}
\newcommand{\ket}[1]{| #1 \rangle}

% Types
\newcommand{\tnat}{\mathsf{Nat}}
\newcommand{\tbool}{\mathsf{Bool}}
\newcommand{\tint}{\mathsf{Int}}
\newcommand{\tposreal}{\mathsf{PosReal}}
\newcommand{\treal}{\mathsf{Real}}
\newcommand{\tunit}{\mathsf{Unit}}
\newcommand{\tprod}{\mathsf{Prod}}
\newcommand{\tP}[2]{P \ #1 \ #2}
\newcommand{\unitgrade}{[]}
\newcommand{\rarrow}{\Rightarrow}

% Type system
\newcommand{\dtyping}[3]{#1 \vdash^d #2: #3}
\newcommand{\ptyping}[4]{#1 \mid #2 \vdash^p #3: #4}
\newcommand{\rtyping}[4]{#1 \vdash #2: #3 \rarrow #4}

\newcommand{\ctxequiv}{\simeq_{\mathsf{ctx}}}
\newcommand{\Values}{\mathsf{Values}}

% DenSem
% \newcommand{\sem}[1]{\llbracket #1 \rrbracket}
\newcommand{\subfin}{\subseteq_{\fin}}
\newcommand{\inj}{\mathbf{inj}}
\newcommand{\detsem}[1]{[#1]}
\newcommand{\List}{\mathrm{List}}
\newcommand{\sIntervention}{\mathrm{Intv}}
\newcommand{\sP}[2]{\mathbf{P} \ #1 \ #2}
\newcommand{\sPP}{\mathbf{P}}
\newcommand{\sVal}{\mathbf{Val}}
\newcommand{\gradecatsymb}{\mathbb{G}}
\newcommand{\gradecat}[1]{\gradecatsymb_{#1}}
\newcommand{\sdist}{\mathrm{dist}}
\newcommand{\sintervene}{\mathrm{intervene}}

\DeclareMathOperator*{\dom}{dom}

% New language
\newcommand{\sCaus}{\mathrm{Caus}}
\newcommand{\Caus}[2]{\sCaus \ #1 \ #2}
% \newcommand{\Maybe}{\mathsf{Maybe}}
\newcommand{\opnew}{\mathsf{new}}
\newcommand{\Caustic}{\mathsf{\lambda_C}}
\newcommand{\cdtyping}[3]{#1 \vdash^d_{\Caustic} #2: #3}
\newcommand{\crtyping}[4]{#1 \vdash^r_{\Caustic} #2: #3 \rarrow #4}
\newcommand{\cptyping}[4]{#1 \vdash^{p}_{\Caustic} #2 :#3 \& #4}
\newcommand{\cptypingreordered}[4]{\cptyping{#2}{#3}{#4}{#1}}

% Compilation and Logical Relation
\DeclareMathOperator*{\len}{dom}

\newcommand{\Maybe}[1]{\mathrm{Maybe}({#1})}
\newcommand{\sem}[1]{\llbracket{#1}\rrbracket}
\newcommand{\Just}[1]{\mathrm{Just}({#1})}
\newcommand{\Nothing}{\mathrm{Nothing}}

% Dario's Macros
\usepackage{bm}
\renewcommand{\vec}{\bm}

\newcommand{\todo}[1]{\textcolor{blue}{[TODO: #1]}}
\newcommand{\prompt}[1]{\textcolor{purple}{[PROMPT: #1]}}

\newcommand{\cat}[1]{\mathbb{#1}}
\newcommand{\catname}[1]{\mathbf{#1}}

\newcommand{\set}{\mathbf{set}}
\newcommand{\Set}{\catname{Set}}

\newcommand{\concat}{+\!\!+}

\newcommand{\defeq}{\coloneq}


\newcommand{\C}{\cat C}
\newcommand{\G}{\mathcal G}

\newcommand{\op}{\mathrm{op}}


\newcommand{\thought}[3]{\marginpar{\tiny\itshape\raggedright\textcolor{#2}{\textbf{#1:} #3}}\ignorespaces} \marginparwidth=1.4cm \marginparsep=0cm
\newcommand{\jeremy}[1]{\thought{Jeremy}{green!20!black}{#1}}

\newcommand{\IntvInstr}{\mathsf{IntvInstr}}
\newcommand{\oplookup}{\mathsf{lookup}}
\newcommand{\oprestrict}{\mathsf{restrict}}
\newcommand{\FinSet}{\catname{FinSet}}
\newcommand{\Bij}{\catname{Bij}}
\newcommand{\Type}{\mathbf{Type}}

\newcommand{\vpure}{\mathbf{pure}}
\newcommand{\vapp}{\mathbin{\langle * \rangle}}
% \DeclareMathOperator*{\vapp}{\langle * \rangle}
\newcommand{\fmap}{\mathbin{\langle \$ \rangle}}

\newcommand{\colorone}[1]{\textcolor{purple}{#1}}
\newcommand{\colortwo}[1]{\textcolor{blue}{#1}}
\newcommand{\ParIndex}{\mathsf{ParIndex}}
\newcommand{\Env}{\mathbf{Env}}

\newcommand{\inl}{\mathbf{inl}}
\newcommand{\inr}{\mathbf{inr}}

\newcommand{\note}{\textcolor{purple}}
\newcommand{\cur}{\mathbf{cur}}
\newcommand{\uncur}{\mathbf{uncur}}

\newcommand{\intvPt}{\mathsf{intvPt}}
\newcommand{\Endo}{\mathcal{E}}
\newcommand{\sEndo}{\mathsf{Endo}}

\newcommand{\arrow}{\rightsquigarrow}
\newcommand{\arrarr}{\mathsf{arr}}
\newcommand{\arrcomp}{\ggg}
\newcommand{\arrfirst}{\mathsf{first}}


\newcommand{\DarioMonad}{\mathsf{T}_{\mathsf{CBN}}}
\newcommand{\oprun}{\mathsf{run}}
\newcommand{\dolet}{\mathbf{let}}
\newcommand{\doletexpr}[3]{\dolet \ #1 \leftarrow #2 \ \doin \ #3}
\newcommand{\doletexprParam}[4]{\dolet_{#1} \ #2 \leftarrow #3 \ \doin \ #4}
\newcommand{\doin}{\mathbf{in}}

\newcommand{\typing}[3]{#1 \vdash #2: #3}
\newcommand{\wtyping}[4]{#1 \vdash_{#2} #3: #4}
\newcommand{\polytyping}[3]{#1 \vdash_{[-]} #2: #3}
\newcommand{\TMon}[2]{T_{#1}\left(#2 \right)}

\newcommand{\catIntv}{\mathbf{IntvStruct}}

\newcommand{\bOne}{\mathds{1}}
\newcommand{\newEndo}{\mathsf{endo}}

\newcommand{\Meas}{\mathsf{Meas}}
\newcommand{\perWorld}{\mathsf{perWorld}}
\newcommand{\valueAt}{\mathsf{valueAt}}

\newcommand{\PSh}{\mathbf{PSh}}
\newcommand{\PShBij}{\PSh(\Bij)}

\newcommand{\LangName}{\mathsf{\lambda_C}}
\newcommand{\World}{\mathsf{World}}

\newcommand{\dwtyping}[4]{#1 \mid #2 \vdash^d #3: #4}
\newcommand{\pwtyping}[5]{#1 \mid #2 \mid #3 \vdash^p #4: #5}
% \newcommand{\rtyping}[4]{#1 \vdash #2: #3 \rarrow #4}

\newcommand{\Any}{\mathsf{Any}}
\newcommand{\intvConstr}{\mathsf{intvConstr}}
\newcommand{\twin}{\mathsf{twin}}

\newcommand{\Functions}{\mathsf{Functions}}
\newcommand{\Distributions}{\mathsf{Distributions}}
\newcommand{\Constants}{\mathsf{Constants}}

\newcommand{\rone}[1]{{#1}^{(1)}}
\newcommand{\rtwo}[1]{{#1}^{(2)}}
\newcommand{\rthree}[1]{{#1}^{(3)}}

\newcommand{\Lone}{L_1}
\newcommand{\Ltwo}{L_2}
\newcommand{\Lthree}{L_3}

\newcommand{\Sbs}{\catname{Sbs}}
\newcommand{\Qbs}{\catname{Qbs}}

\newcommand{\typingone}[4]{#1 \vdash_{1}^{#4} #2: #3}
\newcommand{\etypingone}[3]{\typingone{#1}{#2}{#3}{\mathbf{e}}}
\newcommand{\ttypingone}[3]{\typingone{#1}{#2}{#3}{\mathbf{t}}}
\newcommand{\qtypingone}[3]{\typingone{#1}{#2}{#3}{\mathbf{q}}}

\newcommand{\typingtwo}[5]{#1 \mid #2 \vdash_{2}^{#5} #3: #4}
\newcommand{\etypingtwo}[4]{\typingtwo{#1}{#2}{#3}{#4}{\mathbf{e}}}
\newcommand{\ttypingtwo}[4]{\typingtwo{#1}{#2}{#3}{#4}{\mathbf{t}}}
% \newcommand{\qtypingtwo}[4]{\typingtwo{#1}{#2}{#3}{#4}{q}}

\newcommand{\etypingthree}[4]{#1 \mid #2 \vdash_{3}^{\mathbf{e}} #3: #4}
\newcommand{\ttypingthree}[6]{#1; #2 \mid #3; #4 \vdash_{3}^{\mathbf{t}} #5: #6}
% \newcommand{\qtypingtwo}[4]{\typingtwo{#1}{#2}{#3}{#4}{q}}
\newcommand{\Gammaexo}{\Gamma_{\mathsf{exo}}}
\newcommand{\Gammaendo}{\Gamma_{\mathsf{endo}}}
\newcommand{\Deltaexo}{\Delta_{\mathsf{exo}}}
\newcommand{\Deltaendo}{\Delta_{\mathsf{endo}}}

\newcommand{\typingHO}[3]{#1 \vdash_{\mathsf{HO}} #2: #3}


\newcommand{\scoringquery}[4]{\letbind{#1}{#3}{\opscore(#2)};\ \return{#4}}

\newcommand{\xletbind}[3]{\mathbf{let \ exo} \ #1 = #2 \ \mathbf{in} \ #3}
\newcommand{\nletbind}[3]{\mathbf{let \ endo} \ #1 = #2 \ \mathbf{in} \ #3}


\newcommand{\opdefault}{\mathbf{default}}
\newcommand{\opcondition}{\mathsf{condition}}
\newcommand{\opapplyIntv}{\mathsf{applyIntv}}
\newcommand{\forget}{\mathbf{forget}}
\newcommand{\semSCM}[1]{\sem{#1}_{\mathsf{SCM}}}

% Section 4
\newcommand{\IntvConstrName}{\mathbf{IntvC}}
\newcommand{\semone}[1]{\sem{#1}_{\Lone}}
\newcommand{\semtwo}[1]{\sem{#1}_{\Ltwo}}
\newcommand{\semthree}[1]{\sem{#1}_{\Lthree}}
\newcommand{\semho}[1]{\sem{#1}_{\LHO}}
\newcommand{\liftedEndo}{\widehat{\Endo}}
\newcommand{\sliftedEndo}{\sEndo}
\newcommand{\lift}{\mathsf{lift}}
\newcommand{\liftEndo}{\mathsf{endo}}
\newcommand{\catp}{\mathbf{P}}
\newcommand{\App}{\mathsf{App}}
\newcommand{\intvLabel}{\mathsf{intvLabel}}
\newcommand{\ulBij}{\underline{\Bij}}

% Section 5
\newcommand{\LHO}{L_{\mathsf{HO}}}
\newcommand{\LthreeHOcompile}[1]{\llparenthesis #1 \rrparenthesis}
\newcommand{\IntvMonadSymb}{\mathcal{I}}
\newcommand{\IntvMonad}[2]{\IntvMonadSymb_{#1}\left(#2\right)}
\newcommand{\Twinned}{\mathcal{T}}
% \newcommand{\TwinnedIntvT}[1]{\Twinned_{#1} \IntvMonadSymb}
\newcommand{\sTwinned}{\mathsf{Twinned}}
\newcommand{\sIntvMonadSymb}{\mathsf{IntvT}}
\newcommand{\sIntvMonad}[2]{\sIntvMonadSymb_{#1}\left(#2\right)}
\newcommand{\opnewEndo}{\mathsf{newEndo}}
\newcommand{\unzip}{\mathsf{unzip}}
\newcommand{\zip}{\mathsf{zip}}

% commenting
\colorlet{jyColour}{blue}

\makeatletter
\patchcmd{\@@addmarginpar}{\ifodd\c@@page}{\ifodd\c@@page\@@tempcnta\m@@ne}{}{}
\makeatother
\reversemarginpar
\newcommand{\sidenote}[3]{\marginpar{\tiny\itshape\raggedright\textcolor{#2}{\textbf{#1:} #3}}\ignorespaces} \marginparwidth=1.4cm \marginparsep=0cm
\newcommand{\jy}[1]{\sidenote{jy}{jyColour}{#1}}


% SCMs
\DeclarePairedDelimiterX{\parensbar}[1]
{(\!\delimsize\vert}
{\delimsize\vert\!)}
{#1}
\newcommand{\compileSCM}[1]{\ensuremath{\parensbar{#1}}}
\newcommand{\scmIntervene}[2]{\ensuremath \mathsf{do}_{#1, #2}}
\newcommand{\scmTwin}[1]{\ensuremath \mathsf{twin}_{#1}}
\newcommand{\IntClass}{\ensuremath{\mathsf{Intv}}}
\newcommand{\DefaultClass}{\ensuremath{\mathsf{Default}}}

\newcommand{\commentcolor}{\textcolor{codegreen}}
\newcommand{\highlightcolor}{\textcolor{blue}}
\newcommand{\SCM}{\mathcal{M}}

\usepackage{xcolor}
\usepackage{booktabs,array,tabularx}
\usepackage{listings}
\lstset{basicstyle=\ttfamily\small,mathescape=true,columns=fullflexible,keepspaces=true}

\addbibresource{refs.bib}
\addbibresource{refs-quantum.bib}
\addbibresource{refs-new.bib}
\setcounter{biburlnumpenalty}{100}\setcounter{biburlucpenalty}{100}\setcounter{biburllcpenalty}{100}

\newcommand{\fillin}[1]{{\color{orange!80!black}[FILL: #1]}}
\newcommand{\iotq}{\mathsf{IO}\otimes\mathsf{Q}}
\newcommand{\CPinf}{\mathbf{CP}^{\infty}}
\newcommand{\CP}{\mathbf{CP}}
\providecommand{\LHO}{L_{\mathsf{HO}}}
\newcommand{\Lk}[1]{L_{#1}}
\newcommand{\CausC}{\mathbf{Caus}}
\newcommand{\FinStoch}{\mathbf{FinStoch}}
\newcommand{\Mat}{\mathbf{Mat}(\mathbb{R}_{\geq 0})}
\newcommand{\WP}[1]{WP\,#1}
\newcommand{\opin}{\mathsf{in}}
\newcommand{\opout}{\mathsf{out}}
\providecommand{\opnew}{\mathsf{new}}
\newcommand{\opmeasure}{\mathsf{measure}}
\newcommand{\applyU}[1]{\mathsf{apply}_{#1}}
\newcommand{\dok}{\mathsf{do}}
\newtheorem{conj}{Conjecture}
\newtheorem{hyp}{Hypothesis}


\title{Semantics of Causality:\\ Interventional, Signalling and Factorisation Structure\\ in Probabilistic and Quantum Programs\\[1ex]
\large Qualifying Dissertation for the Transfer of Status}
\author{Theo Wang\\ \small Department of Computer Science, University of Oxford\\ \small Supervisor: Sam Staton}
\date{Michaelmas Term 2026}

\begin{document}
\maketitle

\begin{abstract}
\noindent
The word \emph{causality} names several different things. In causal inference it means invariance under \emph{interventions}; in quantum foundations it means the absence or presence of \emph{signalling}; in probabilistic modelling it means a \emph{factorisation} of a joint distribution along a graph. All three are extensional characterisations of an intensional datum that programs carry explicitly: the dataflow of a probabilistic program, the order of communication events of a communicating program, the graph of a causal model. Two lines of work precede this proposal. The first, carried out during my first year with coauthors, is a hierarchy of causal probabilistic programming languages in which interventional causal structure is characterised by contextual equivalence. The second, begun in my MSc dissertation and continued this year, is an algebraic theory of classically controlled quantum communication in which the order of I/O events is first-class semantic data. The proposed thesis is to give the second setting its signalling reading, to relate the two by a functorial bridge on their common finite classical fragment, and, as a stretch goal, to axiomatise a categorical framework in which all three notions of causality are structure on one category. Part I reviews the literature the thesis builds on and positions the two neighbouring fields it reaches into; Part II states the research plan and is self-contained.
\end{abstract}

\tableofcontents

% =====================================================================
\part{Literature Review}
% =====================================================================

\section{Introduction}\label{sec:intro}

\paragraph{Three notions of causality.}
When a computer scientist says that one part of a system \emph{causes} another, they may mean at least three distinct things. In Pearl's theory of causal inference \cite{pearl2009causality}, $X$ has a causal effect on $Y$ if \emph{intervening} on $X$---forcing it to a value, regardless of how it would otherwise have been produced---changes the distribution of $Y$. In quantum foundations and in the theory of higher-order quantum processes \cite{kissinger2019categorical,oreshkov2012noncausal}, the causal structure of a process is given by its \emph{signalling} constraints, which inputs can influence which outputs, together with the further question of whether it is compatible with a definite causal order between its parts. In probabilistic modelling, a causal structure is a directed acyclic graph along which a joint distribution \emph{factorises}, and the content of the structure is the family of conditional independences it entails \cite{pearl2009causality,fritz2023dseparation}.

Three further usages should be named so as to be set aside. In categorical quantum mechanics, a process is called \emph{causal} when it preserves the discarding map \cite{PQP_Coecke_Kissinger_2017}. This is the causality axiom of Chiribella, D'Ariano and Perinotti \cite{chiribella2010purification}, glossed as ``no signalling from the future'', and Coecke \cite{coecke2014terminality} shows that it implies the no-signalling principle for processes composed along a causal structure; it is the base axiom on which signalling structure is built, not a notion of causality between parts. Fritz \cite{fritz2020synthetic} calls a Markov category \emph{causal} when the output of a morphism cannot depend on a choice made downstream of it; this is again a property of the ambient category. In artificial intelligence and in verification, \emph{actual causality} \cite{halpern2016actual} asks which particular event caused which particular outcome in a single run; that question is not treated here. Nor are experimental tests of indefinite causal order, causal discovery, or identifiability, each of which is mentioned below only where a work package touches it.

That the same word carries these meanings is not an accident, but neither is it a theorem. Each notion is \emph{extensional}: it characterises causal structure by what an observer may do to the system---intervene, send signals, condition. What such observers probe is \emph{intensional}: the dataflow of a program, the temporal order of its communication events, the graph of a model. Programs carry this intensional structure explicitly in their syntax. The viewpoint of this thesis is therefore that causality is a property of \emph{program semantics}, and that the right way to relate the three notions is to relate, for each one, the intensional structure in the syntax to the extensional characterisation in the semantics by a correspondence theorem.

\paragraph{Where the two routes stand.}
Two lines of work precede this proposal, and the review is organised around them. On the interventional side, my coauthors and I built a hierarchy of first-order probabilistic languages in which Pearl's interventions and counterfactual \emph{twinning} are program transformations, defined program equivalences contextually, and gave a graded denotational semantics that is fully abstract for concrete models (\S\ref{sec:intv}). This route has the complete picture for one notion of causality: the intensional dataflow of a program is characterised extensionally by equivalence under all interventional contexts. On the communication side, my MSc dissertation gave an algebraic theory $\iotq$ of classically controlled quantum communication, an operational model in which programs exchange qubits with a quantum stream, and a denotational model given by a free I/O monad on $\CPinf$, with the denotational model adequate and fully abstract for the operational one; this year's work, with my supervisor, is its conference write-up and a draft proof of completeness (\S\ref{sec:order}). Here the temporal order of communication events is first-class semantic data: the denotation is indexed by I/O traces. This route does \emph{not} yet have an extensional causal reading of that order; giving it one, as signalling structure in the sense of higher-order quantum processes, is part of the research plan rather than a result.

A single example shows why these are two faces of one question. In the interventional language, a program reads an \emph{intervention label} $\ell$ through $\mathsf{applyIntv}\ \ell\ t$, which yields the label's value when it is set and $t$ otherwise; if a second variable $z$ does not depend on the result, every intervention at $\ell$ leaves the distribution of $z$ unchanged, and the contextual equivalence makes this an equation between programs. Causal independence is \emph{tested by contexts}. In the theory $\iotq$, the commutation axiom $\applyU{U}(\vec a,\ \vec a.\ \opin(b.\ x(\vec a, b))) = \opin(b.\ \applyU{U}(\vec a,\ \vec a.\ x(\vec a, b)))$ says that a local operation and a communication event on different qubits may be reordered, and the quantum-stream semantics validates it; by contrast $\opin$ and $\opout$ do not commute with one another, and that is exactly why the trace index is a total order. Causal independence is \emph{an equation validated by an observer}. Both statements say that no information flows from one place to another, so that reordering or intervening is unobservable; in programming-language terms both are \emph{noninterference} properties \cite{goguen1982noninterference}. Whether this coincidence is an instance of a general correspondence is the question of Part II.

\paragraph{Scope of the review.}
Sections \ref{sec:foundations}--\ref{sec:order-wherewestand} treat the literature on which the thesis builds: the semantics of effects, probabilistic programs and quantum programs; causal models, their equivalences and their diagrammatic presentations; causal probabilistic programming; and quantum process calculi. Sections \ref{sec:concurrency} and \ref{sec:hoqc} position two neighbouring fields that the research plan reaches into, the semantics of concurrency and higher-order quantum processes, at the level of their main definitions and results; their detailed survey is a milestone of that plan (Part II, \S\ref{sec:plan}). Section \ref{sec:synthesis} draws the three notions together and identifies the gap the thesis addresses.

% ---------------------------------------------------------------------
\section{Semantic foundations shared by both routes}\label{sec:foundations}

\subsection{Monads, algebraic theories, and their combinations}

Moggi's observation that notions of computation are monads \cite{MOGGI199155,moggi1991monadicmetalanguage} and Plotkin and Power's refinement that most of them are presented by algebraic theories of operations and equations \cite{plotkin_power_2002,PlotkinPower04CompEffectsOverview} are the common foundation of both routes. An algebraic theory determines a monad, and conversely a finitary monad is a Lawvere theory \cite{lawveretheories_monads_HYLAND2007437}; the equational presentation is what makes program equivalence a matter of derivability. Hyland, Plotkin and Power \cite{sums_tensors_HYLAND200670} showed that the standard ways of combining theories have computational meaning: the \emph{sum} adjoins the operations of two theories with no interaction, so that the order of effects from the two theories is observable, while the \emph{tensor} additionally imposes that operations of one theory commute with those of the other, so that their relative order is unobservable. In the premonoidal semantics of effects \cite{premonoidal_cat_POWER_ROBINSON_1997,PowerRobinson97premonoidal,fuhrmann1999premonoidal}, the morphisms that commute with everything are the \emph{central} ones; centrality is the established algebraic form of ``cannot be observed to happen before or after''. Handlers \cite{bauer2019algebraicalgebraiceffectshandlers} give the programming counterpart of models; comodels and runners \cite{model_comodel_PLOTKIN2008295,runners_10.1007/978-3-030-44914-8_2} describe the environments in which effectful programs execute.

Ordinary algebraic theories cannot express operations that create and consume \emph{resources}, such as qubits, which by no-cloning must be treated linearly. Staton's \emph{parameterised algebraic theories} (PATs) \cite{statonlics13,statonpopl15} extend the framework with parameters: operations consume parameters, classically choose a branch, and bind fresh parameters in each branch. Staton used them to give a fully complete equational theory of qubit quantum computation \cite{statonpopl15}; Stark's free model of $\pi$-calculus name creation \cite{stark:freamp} is the earlier instance of the same commutation discipline, where communication primitives commute with independent local operations. Kammar, Liell-Cock, Lindley, Matache and Staton \cite{kammar2026threads} have given a complete equational axiomatisation of dynamic threads (fork, wait, stop, atomic actions) as a parameterised algebraic theory, whose closed terms correspond exactly to pomsets; I return to this in \S\ref{sec:concurrency}, since it is the nearest neighbour of my own completeness result.

\emph{Where we stand.} My MSc dissertation \cite{wang_algebraic_2024} shows that sums and tensors exist for the enriched Lawvere theories presented by PATs, by explicit syntactic construction, which is what licenses the definition of $\iotq$ as a tensor in \S\ref{sec:order}.

\emph{What this gives us.} The tensor, and centrality more generally, are already primitive statements about causal independence. In $\iotq$ the tensor makes communication commute with local operations on other qubits, while $\opin$ and $\opout$ do not commute with each other. I treat this as an observation to be made precise, not as a result: Part II asks (\WP{3}) whether centrality in an algebraic theory is in general the algebraic shadow of no-signalling.

\subsection{Semantics of probabilistic programs}

Probabilistic programming gives a compositional syntax for statistical models and performs Bayesian inference on them compositionally \cite{vandemeent2021introductionprobabilisticprogramming}. Its denotational semantics is monadic: the Giry monad on measurable spaces \cite{kock72,semantics_ppl_staton2016}, and, for higher-order programs, the probability monad on quasi-Borel spaces \cite{Heunen_2017}. Implementations such as LazyPPL \cite{lazyppl}, MonadBayes \cite{monadbayes1,monadbayes2} and GenJAX \cite{becker26genjax} realise this monadic structure directly. Full abstraction is the recurring theorem shape for such languages, the standard reference point being probabilistic coherence spaces for probabilistic PCF \cite{ehrhard2014pcoh}.

The synthetic approach of \emph{Markov categories} \cite{fritz2020synthetic,fritz2019probability} axiomatises the categorical structure needed to talk about probability without measure theory: a symmetric monoidal category with a commutative comonoid on every object, where discarding is natural (the unit is terminal) but copying need not be. In Markov categories with conditionals, disintegration and Bayesian inversion are expressed synthetically \cite{cho2019disintegration,fritz2020synthetic}, and Fritz and Klingler \cite{fritz2023dseparation} prove the soundness of d-separation, and the equivalence of the local and global Markov properties, synthetically; completeness of d-separation in the sense of Geiger, Verma and Pearl remains the classical result. This is the \emph{factorisation} face of causality: a Bayesian network is a string diagram in a Markov category, and the conditional independences it entails are read off the diagram.

The second piece of structure both routes rely on is \emph{grading}. Graded monads \cite{katsumata2014_graded_strong_monad,orchard2020unifying,fujii2016towards,dorsch2019graded} and, more generally, locally graded categories \cite{levy2019locallygraded,sarkis2025string} track an index---an effect annotation, a trace, a resource---through composition. Milius, Pattinson and Schröder \cite{milius2015gradedtrace} use graded monads for trace semantics; Liell-Cock and Staton's $\mathbf{ImP}$ construction \cite{liell2025compositional} uses grading to make imprecise probability compositional. The interventional semantics of \S\ref{sec:intv-wherewestand} is built on an analogous construction in which the grade records which interventions are available.

\emph{What this gives us and what it lacks.} Markov categories give factorisation but no primitive notion of intervention; the diagrammatic causal-model literature (\S\ref{sec:diagrammatic}) adds it by surgery on diagrams. Neither offers a notion of program \emph{context}, which is what contextual equivalence needs.

\subsection{Semantics of quantum programs}

Selinger's QPL \cite{QPL_SELINGER_2004} established the denotational template for quantum programming languages: types denote finite-dimensional spaces of density matrices (or, more generally, $C^*$-algebras \cite{SemOperatorAlgebra}), and programs denote completely positive maps. Categorical quantum mechanics \cite{abramsky2007categoricalsemanticsquantumprotocols,heunen_vicary_book,PQP_Coecke_Kissinger_2017} gave this a graphical calculus through dagger compact closed categories and the CPM construction \cite{selinger_dagger_compact_closed_cat_cp_maps}, and the ZX-family of calculi axiomatise fragments of it; Staton's PAT $\mathsf{Q}$ \cite{statonpopl15} and Carette et al.'s extension of $\Pi$ \cite{Quantum_PI_Carette_2024} give complete equational theories of qubit programs and of unitary circuits respectively. Languages such as QWIRE, Quipper and Silq \cite{qwire_10.1145/3009837.3009894,quipper_DBLP:journals/corr/abs-1304-3390,silq} work in the circuit model with classical control. For languages with infinitary features one needs the completion of $\CP$ under infinite biproducts, $\CPinf$ \cite{PaganiSV13,D_Completion_ZHAO20102167}, which is the base category of the communication semantics. Quantum monads \cite{sati2023quantummonadology,QIO_Altenkirch_Green_2009} and quantum game semantics \cite{clairambault2019quantumgames} are two further lines of work on the structure of quantum computation as an effect; the latter uses event structures and is the closest existing semantics to the one proposed in \WP{2a}.

\emph{What this gives us.} The axiom that singles out the deterministic (discard-preserving) processes in categorical quantum mechanics \cite{PQP_Coecke_Kissinger_2017,chiribella2010purification} is the base condition from which the signalling structure of \S\ref{sec:hoqc} is built; the $\CausC$ construction requires compact closure and further ``precausality'' axioms on top of it, and, by design, its base category also contains the non-deterministic, trace-decreasing processes.

% ---------------------------------------------------------------------
\section{Causality as intervention}\label{sec:intv}

\subsection{Structural causal models and Pearl's hierarchy}

Pearl's structural causal models (SCMs) \cite{pearl2009causality,Pearl2018_book_of_why} are systems of equations $X_i = f_i(\mathrm{Pa}_i, U_i)$ in which each endogenous variable is a deterministic function of its parents and of an exogenous noise variable; a prior on the noise makes an SCM a probabilistic model, and the equations make it more than a distribution. The surplus is organised by the \emph{causal hierarchy} \cite{bareinboim:etal20}: rung 1 asks observational questions; rung 2 asks \emph{interventional} questions, answered by replacing the equation for $X$ with a constant (the $\dok$ operator); rung 3 asks \emph{counterfactual} questions about a particular run, answered by abduction, action and prediction, implemented graphically by \emph{twinning}, which duplicates the model into worlds that share the noise \cite{Balke1994}. Both operations are \emph{syntactic transformations of the model}. The Neyman--Rubin potential-outcomes framework \cite{rubin1974} is the other dominant formalism; its correspondence with SCMs is well understood \cite{ibeling2023comparing}. Richardson and Robins' single-world intervention graphs \cite{richardson2013swig} present an intervention by splitting a node into its natural and its fixed value, which is the graphical counterpart of the intervention labels of our languages; they are deliberately single-world, and cross-world quantities such as natural direct effects need further assumptions there. Interventions need not be perfect: the literature on stochastic and soft interventions \cite{dawid2002_stoch_intv,didelez2006_stoch_mech_intv,Eberhardt_Scheines_2007,levis2024stochasticinterventionssensitivityanalysis} replaces an equation by another mechanism rather than a constant.

Throughout, I work with fully specified models and Bayesian inference: mechanisms and priors are given, and queries are answered by inference. \emph{Identifiability} (when an interventional quantity is computable from observational data, the subject of the do-calculus \cite{JMLR:v9:shpitser08a,shpitser2012counterfactualstested}) and \emph{causal discovery} \cite{spirtes2000causation,peters_janzing_schoelkopf_2017_book} are outside the scope of this thesis.

\emph{What this gives us.} The causal structure of an SCM is its dataflow: which variables each equation reads. Interventions probe it. The semantic question is when two models have the same causal structure, and the answer must be extensional.

\subsection{Equivalences of causal models and logical accounts}

Bongers, Forré, Peters and Mooij \cite{Bongers_2021} give the principal account of equivalence for SCMs at all three rungs: two models are observationally, interventionally or counterfactually equivalent when they agree on the corresponding class of queries, with the counterfactual case defined through the two-world twinned model. Their treatment also covers cyclic SCMs, through solution functions. At rung 2, several further notions exist: Pearl's causal Bayesian networks are families of interventional distributions compatible with a graph; \emph{causal spaces} \cite{park2024measuretheoreticaxiomatisationcausality,buchholz2024products} axiomatise this measure-theoretically; and \emph{causal abstraction} \cite{rubenstein2017consistency,beckers2019abstracting,rischel2021abstraction,lorenztull2026abstraction,englberger2025causalabstractionscategoricallyunified} studies coarsening maps between models that commute with interventions. At the level of graphs rather than models, Markov equivalence \cite{verma1990equivalence} is what discovery algorithms can recover. On the logical side, Ibeling and Icard \cite{ibeling2020hierarchy} characterise the three rungs by three probabilistic logics of strictly increasing expressivity whose semantics is defined on SCMs through Pearl's operations, and Icard has proposed probabilistic Turing machines as the computational counterpart of SCMs \cite{icard2017programs,ibeling2019openuniverse}.

\emph{What this gives us and what it lacks.} These are all extensional characterisations of interventional structure, and they are the yardstick against which our contextual equivalences are measured (\S\ref{sec:intv-wherewestand}). None of them is formulated as a contextual equivalence of programs, and none treats higher-order models.

\subsection{Diagrammatic and categorical causal models}\label{sec:diagrammatic}

Fong's \emph{causal theories} \cite{fong2012causaltheories} present a Bayesian network as a symmetric monoidal category freely generated by its graph, with models as functors into stochastic maps. Jacobs, Kissinger and Zanasi \cite{jacobs2019causal} present interventions as \emph{string-diagram surgery} in this setting, and introduce \emph{comb disintegration}: the use of combs, processes with a hole, to express the classical causal-inference computations, which is the first appearance of the comb notion of \S\ref{sec:hoqc} in classical causal inference. Lorenz and Tull \cite{lorenz2023causalmodelsstringdiagrams} develop causal models in string diagrams systematically, and a later line presents causal abstraction as natural transformations between such models \cite{lorenztull2026abstraction,englberger2025causalabstractionscategoricallyunified}. Bayesian inference in this style is handled by \emph{normalisation boxes}, which renormalise locally inside a diagram \cite{jacobsszelesstein2025inference}, and Agrawal et al.\ \cite{agrawal2026equational} obtain equational principles for causal Bayesian networks from them. Zhang et al.\ \cite{zhang2026semanticseparation} characterise d-separation as a noninterference-style separation property, which makes explicit a connection implicit in all of this work and in \S\ref{sec:intro}: invariance under interventions is a noninterference property \cite{goguen1982noninterference}.

\emph{What this gives us and what it lacks.} This is the closest existing bridge between the factorisation face and the intervention face of causality, and it is compositional in the functorial sense. Its equational theories are, however, stated per graph, and it has no notion of program context or of higher-order models. It is where the unifying work package (\WP{3}) starts, and comb disintegration is the prior work that \WP{3a} must be positioned against.

\subsection{Causal probabilistic programming languages}

Causal probabilistic programming aims to make intervention and twinning primitives of an expressive probabilistic language. Libraries such as DoWhy \cite{pywhy,blobaum2024dowhygcm} and the $\dok$ operators of Pyro and PyMC \cite{bingham2018pyro,abrilpla2023pymc} manipulate models as graphs or computation graphs and stop at rung 2 (DoWhy's GCM module reaches rung 3 on fitted mechanisms). Witty et al.\ \cite{witty2019synthesis,witty2023bayesian} treat an intervention as a source transformation that replaces subterms, which is powerful but non-compositional; Brulé's Whittemore \cite{brule2018whittemore} is an early embedded causal DSL.

Four systems reach rung 3 natively. OmegaC \cite{tavares2021language} obtains interventions by a lazy evaluation strategy that retroactively modifies closures, reminiscent of dynamic binding, and counterfactuals by making every program a deterministic function of explicit noise. ChiRho \cite{chirho} implements both as effect handlers over Pyro, with a multi-world value type whose relation to Pearl's twinning is not immediate. MultiVerse \cite{perov2020multiverse} uses trace inspection and re-execution. Kiesel et al.\ \cite{kiesel2023whatif} answer counterfactual queries on ProbLog programs by a transformation proved correct against Pearl's semantics.

\emph{What this gives us and what it lacks.} These systems show that the primitives are implementable at higher order, by non-standard evaluation techniques. What they lack is a denotational and equational foundation: a semantics in which intervention is a mathematical operation rather than an evaluation strategy, and an account of when two causal programs are equivalent. That is what the first route supplies.

\subsection{Where we stand on the interventional route}\label{sec:intv-wherewestand}

Our POPL submission \fillin{status: under revision for POPL 2027?; coauthors} (an early version appeared at the LAFI workshop \cite{lafi-wang}; artifact \cite{ARTIFACT_wang_2026_23207221}) reconstructs Pearl's hierarchy as a hierarchy of three first-order probabilistic languages. $\Lk{1}$ is a standard first-order probabilistic language with conditioning by a final score. $\Lk{2}$ exposes hooks in the dataflow as typed free variables, the \emph{intervention labels}, read through $\mathsf{applyIntv}$; an intervention $\dok(\ell\mapsto v)\{t\}$ is substitution of $v$ for $\ell$. $\Lk{3}$ annotates each binding as exogenous or endogenous, and \emph{twinning} $\mathsf{twin}_n$ is a vectorisation transformation from $\Lk{3}$ to $\Lk{2}$ that duplicates the endogenous bindings into $n$ worlds sharing the exogenous ones. Program equivalence at rung $k$ is \emph{contextual} with respect to rung-1 denotational equality, in Milner's sense: two programs are $\equiv_k$ when they are denotationally equal in every rung-$k$ context, and varying the admissible interventions and world counts gives a spectrum of equivalences. For representable acyclic SCMs compiled into $\Lk{3}$, the observational, interventional and counterfactual equivalences of Bongers et al.\ coincide with the equivalences $\equiv_3^{\mathrm{obs}}$, $\equiv_3^{\mathrm{intv}}$ and $\equiv_3^{\text{2-world}}$ of the compiled programs; the full $\equiv_3$ quantifies over all $n$-world twinnings, and whether it is strictly finer than $\equiv_3^{\text{2-world}}$ is open. The denotational semantics interprets interventions by \emph{grading}, in a construction reminiscent of $\mathbf{ImP}$ \cite{liell2025compositional}, and twinning by world-polymorphic vectorisation. It is adequate for the contextual equivalences unconditionally and fully abstract when the underlying model is \emph{concrete}, meaning that values correspond to global elements and global elements separate morphisms; $\mathbf{Set}$, standard Borel spaces and quasi-Borel spaces are concrete, and the construction is parametric in the model. The denotation then motivates a higher-order language $\LHO$ with cross-world interventions, into which $\Lk{3}$ compiles in a semantics-preserving way, and which has a definitional Haskell implementation. \fillin{one sentence on which parts are yours: e.g.\ the contextual-equivalence framework and the full-abstraction proofs}

The limitations relevant to Part II are these: the languages are acyclic; the relation between the contextual equational theory and the normalisation-box equational theories of \S\ref{sec:diagrammatic} is not worked out; the first-order semantics is over normalised probability, while $\LHO$ widens to the unnormalised measure monad to internalise conditioning, so that programs with observations denote sub-normalised maps; and the models are measure-theoretic rather than finite-dimensional, which matters for the signalling face (\S\ref{sec:synthesis}).

% ---------------------------------------------------------------------
\section{Causality as order and signalling in communicating quantum programs}\label{sec:order}

\subsection{Quantum communication and distributed quantum computing}

Quantum communication underlies teleportation \cite{Teleportation_PhysRevLett.70.1895}, key distribution \cite{BB84_BENNETT20147} and, increasingly, proposals for distributed quantum computing that connect small quantum processors into larger ones \cite{DQC_survey_Caleffi_2024}. Programming models for distributed quantum computing are at present circuit-level: they describe which gates are applied where and which qubits are sent when. On the semantic side, Ceragioli, Gadducci, Lomurno and Tedeschi observe that the existing quantum process calculi do not agree on what the observational properties of quantum values should be \cite{lqCCS_10.1145/3632885}. There is, in other words, no agreed account of what a distributed quantum program \emph{means}, nor of which of its behaviours are physically realisable.

\subsection{Quantum process calculi}

Quantum communication has mostly been studied through quantum process calculi---CQP \cite{CQP_gay2004communicatingquantumprocesses}, QPAlg \cite{QPAlg_lalire2004processalgebraicapproachconcurrent}, qCCS \cite{qCCS_ying2010algebraquantumprocesses} and lqCCS \cite{lqCCS_10.1145/3632885}---which extend classical calculi \cite{ccs_10.5555/539036,CSP_10.1145/359576.359585,pi_calc_MILNER19921} with qubit message passing. They are presented operationally, by labelled transition systems over configurations, with program equivalence given by bisimilarity. As the authors of lqCCS show, the proposed bisimilarities differ in subtle ways because of the interaction of concurrency, synchronisation and local nondeterminism with quantum effects; their later coalgebraic account of effect-labelled transition systems \cite{ceragiolicoalgebraic} is an attempt to organise these subtleties. A separate difficulty is well known from the classical theory: parallel composition is not an algebraic operation and cannot be adjoined to an algebraic theory as an operation \cite{stark:freamp,CSP_van_Glabbeek_2010}, although workarounds exist, such as the cooperative threads of Abadi and Plotkin \cite{Cooperative_Threads_DBLP:journals/corr/abs-1009-2405}.

\emph{What this gives us and what it lacks.} Process calculi are the only existing semantic treatment of communicating quantum programs, and they are expressive. They are operational only, their equivalences are not canonical, and they entangle communication with concurrency, which makes it hard to isolate what is specific to \emph{communication}. My work studies communication in isolation, as an algebraic effect, and is in that sense the local sublanguage of these calculi with a denotational semantics.

\subsection{Where we stand on the communication route}\label{sec:order-wherewestand}

My MSc dissertation \cite{wang_algebraic_2024}, and its conference version with my supervisor \fillin{FoSSaCS submission status}, introduce $\iotq$, an algebraic theory of classically controlled qubit communication, as the \emph{tensor} of a standard I/O theory \cite{sums_tensors_HYLAND200670} with Staton's complete theory $\mathsf{Q}$ \cite{statonpopl15}: programs may send and receive qubits depending on measurement outcomes, and communication commutes with local quantum operations. The tensor is argued to be needed, since one equation that holds in the tensor but not in the sum is required to validate a quantum phenomenon involving entanglement across the program--environment boundary (Example B in Part II). The theory is given two models. The operational model lets a program exchange qubits with an environment modelled as a \emph{quantum stream}; a program equivalence $\approx_q$ compatible with quantum theory is defined from it by inspecting final quantum states, and terms modulo $\approx_q$ form a model of $\iotq$, which shows that the tensor does not collapse in the manner of Eckmann--Hilton. Although operational, $\approx_q$ is therefore not a contextual equivalence; whether a contextual equivalence observing only traces, continuations and a final measurement coincides with it is a question I return to in Part II. The denotational model is the free I/O monad $T(X)=\bigoplus_{\sigma}\mathrm{qbit}^{|\sigma|}\otimes X$ on $\CPinf$, indexed by I/O traces $\sigma$. The main theorem is that the denotational semantics is adequate and fully abstract for the operational one: $\llbracket c\rrbracket = \llbracket c'\rrbracket$ iff $c\approx_q c'$. What the dissertation left as a conjecture is \emph{completeness}: that $\llbracket t\rrbracket=\llbracket t'\rrbracket$ implies $t=t'$ is derivable in $\iotq$. I now have a draft proof of this (manuscript in preparation, not yet closed), by induction on the I/O depth of terms, using a reorganisation step that normalises every term into an I/O-free head applied to a single input and a single output branch, and a factorisation lemma on completely positive maps that lets semantically equal terms be rewritten to a common prefix. If it closes, the three semantics---algebraic, operational, denotational---coincide.

The observation that sets up Part II is this. The trace $\sigma$ is a \emph{temporal} order of communication events, exposed in the denotation: the $\sigma$-component of a closed term is a completely positive map with $|\sigma|$ open qubit wires, one per event, to be used in the order $\sigma$, and when the term branches on measurements this component is sub-normalised, its trace being the probability of the branches that produce $\sigma$. Whether the family of these components is a \emph{causal} object in the signalling sense of \S\ref{sec:hoqc}, a comb or a tester whose slots are ordered by $\sigma$, is a question the dissertation raised and did not answer. It is stated as Conjecture \ref{conj:b1} in Part II.

\subsection{Order in concurrency theory}\label{sec:concurrency}

Concurrency theory has long modelled causal order explicitly. Lamport's happens-before relation \cite{lamport1978time} is the partial order an observer of a distributed system can reconstruct. Mazurkiewicz traces \cite{diekert1995traces} are words modulo the commutation of independent letters, the cheapest partial-order model and the one that arises when independent channels are added to a sequential program. Pratt's pomsets \cite{pratt1986pomsets}---partially ordered multisets of events, up to isomorphism, with series and parallel composition---generalise words to partially ordered words, and Gischer \cite{gischer1988pomsets} axiomatised the equational theory of series-parallel pomsets. Winskel's event structures \cite{winskel1986eventstructures,nielsen1981petri,winskel1995models} add a conflict relation to the causal order and give a domain-theoretic account of nondeterministic concurrent computation. Concurrent Kleene algebra \cite{hoare2011cka} adds a commutative parallel operator to Kleene algebra, related to sequential composition by an exchange law; Kappé, Brunet, Silva and Zanasi \cite{kappe2018cka} prove that series-parallel pomset languages closed under Gischer's subsumption order are its free model, by reduction to an earlier completeness theorem for bi-Kleene algebra \cite{laurence2014completeness}. Within the algebraic-effects tradition, Dvir, Kammar and Lahav \cite{dvir2022sharedstate} give a monadic semantics of shared-state concurrency from an equational theory, and Kammar et al.\ \cite{kammar2026threads} axiomatise dynamic threads as a parameterised algebraic theory whose closed terms are exactly pomsets, with a completeness theorem. Clairambault, de Visme and Winskel \cite{clairambault2019quantumgames} give an event-structure game semantics for quantum programs.

\emph{Position.} The trace-indexed monad of \S\ref{sec:order-wherewestand} is the degenerate single-channel case of this picture: a single channel forces a total order, so traces are words. Two extensions are distinguished in Part II. Adding independent channels to a \emph{sequential} program makes events on different channels commute, so that traces become Mazurkiewicz traces and the monad a free monad over a partially commutative monoid; this keeps the proof architecture of the single-channel case. Adding \emph{threads} makes traces pomsets, and the theory of Kammar et al.\ is a parameterised algebraic theory with exactly that free model, so ``tensor $\iotq$ with a thread theory'' is a well-posed, and much harder, second step. Completeness results in this literature also go through free models; whether their techniques transfer to the quantum setting, where a factorisation lemma on CP maps did the work, is part of the first reading milestone of the plan. The logics of session types are the other point of contact: BV and its relatives have been used as logics of multiparty protocols, and Simmons conjectures that causal consistency of a wiring of black boxes corresponds to deadlock freedom of the protocol \cite{simmons2024thesis}; a trace-indexed semantics of communicating programs is where such a correspondence could be made precise.

\subsection{Higher-order quantum processes and signalling}\label{sec:hoqc}

Higher-order quantum computation studies transformations of quantum processes rather than of states. Chiribella, D'Ariano and Perinotti's \emph{quantum combs} \cite{quantum_circuit_architecture_Chiribella_2008,chiribella2009networks} are processes with holes in a fixed causal order; their probabilistic counterparts, \emph{testers}, and Gutoski and Watrous's \emph{strategies} \cite{gutoski2007games} describe multi-round interactions with sub-normalised outcomes, and the general theory of supermaps forms a hierarchy of types \cite{bisio2019hoqt}. Oreshkov, Costa and Brukner's \emph{process matrices} \cite{oreshkov2012noncausal} are the most general bipartite process operators compatible with quantum mechanics holding locally in each party's laboratory; the multipartite theory and the notion of \emph{causal separability}, in which the order of later parties may depend classically on earlier outcomes, are due to Araújo et al.\ and Oreshkov and Giarmatzi \cite{araujo2015witnessing,oreshkov2016separable}. Some process matrices are causally non-separable. Two kinds of evidence distinguish them: device-independent violation of a \emph{causal inequality}, which Oreshkov, Costa and Brukner's process exhibits, and device-dependent witnesses of non-separability. The \emph{quantum switch} \cite{chiribella2013switch}, which applies two operations in an order controlled by a qubit in superposition, is causally non-separable \cite{araujo2015witnessing} but produces only causal correlations; photonic experiments implement it \cite{procopio2015experimental}, and whether such implementations realise or simulate indefinite order is debated \cite{oreshkov2019timedelocalized,wechs2023timedelocalised}.

Which indefinite-order processes can be built as circuits is the subject of Wechs, Dourdent, Abbott and Branciard \cite{wechs2021qcqc}, who define nested classes of quantum circuits with fixed order, with \emph{classical} control of causal order (the next operation may depend on classical outcomes of earlier ones) and with \emph{quantum} control of causal order, and characterise each by conditions on process matrices. The switch lies in the quantum-control class and not in the classical-control one; every classical-control process is causally separable; and quantum-control processes generate only causal correlations. That the quantum-control class exhausts the physically realisable processes is a proposal, not a theorem.

Categorically, Kissinger and Uijlen's $\CausC$ construction \cite{kissinger2019categorical} builds, over a \emph{precausal} category---compact closed, with a chosen discarding structure and further axioms, the examples being $\CP$ on finite-dimensional systems and $\Mat$, finite matrices of non-negative reals---a $*$-autonomous category whose types are signalling constraints: first-order processes, non-signalling processes, one-way signalling processes, and so on. The dual of the non-signalling type contains all bipartite process matrices, causally separable or not; separability is a property inside that type, not a type, and Simmons proves that every process matrix is an affine combination of definite-order processes \cite{simmons2024thesis}. Simmons's thesis \cite{simmons2024thesis} is the fullest development of this line and the reference for most of what Part II needs on the signalling side. It rebuilds the construction over \emph{additive} precausal categories (compact closed with products, hence biproducts, and a finite causal basis for every object), which gives $\CausC$ products, coproducts and probabilistic choice as type operators; it adds a non-commutative sequential operator for one-way signalling, under which $\CausC$ models Guglielmi's BV-logic \cite{simmons2022bv}; it defines $n$-ary \emph{graph types} encoding compatibility with an arbitrary DAG of signalling relations; it surveys candidate ways of admitting sub-normalised, partial processes into the construction, effectus-style, and where each falls short; and it introduces \emph{causal logic}, pomset logic with directed axiom links, whose proof nets characterise exactly the causally consistent wirings of black boxes \cite{simmons2024thesis,simmons2024completelogic}, with BV and pomset logic separated by a process-matrix example. The finite-causal-basis axiom excludes infinite-dimensional and measure-theoretic bases, which is one of the two obstructions discussed in \S\ref{sec:synthesis}. Finally, \emph{quantum causal models} transfer Pearl's framework to quantum processes: Costa and Shrapnel \cite{costa2016quantumcausalmodelling} gave the first process-matrix-based definition, Allen et al.\ \cite{allen2017quantumcommoncauses} the common-cause line, and Barrett, Lorenz and Oreshkov \cite{barrett2019quantumcausalmodels} the general framework, in which a node is an input--output pair, the process operator factorises as a product of commuting operators (a quantum Markov condition), interventions are local instruments at nodes, and d-separation in the graph implies conditional independence of classical outcomes. The cyclic extension \cite{barrett2021cyclic} gives a cyclic-causal-model account of some processes without definite order, including the switch. Coherent control of channels \cite{coherent_control_channel_Abbott_2020,resource_theory_comm_Kristjánsson_2020} and, on the language side, the PBS-calculus of Clément and Perdrix \cite{clement2020pbs}, which has the switch as a term, are the relevant prior work on coherent control.

\emph{Position.} Three things follow for this thesis. First, the family of $\sigma$-components of an $\iotq$ denotation has the shape of a tester or strategy whose slots are ordered by the traces; that it \emph{is} one is Conjecture \ref{conj:b1}. Second, signalling types are the extensional face of order in the same way that interventional equivalence is the extensional face of dataflow, and quantum causal models are the one existing place where factorisation, intervention and signalling are all present, which is why they are the first reading milestone. Third, the analogy between parallel composition and indefinite causal order must be stated with care, and the process-matrix picture is the right way to state it: the parties' local programs are combs, and the process matrix is the \emph{network} that wires them together, in other words the scheduler. Parallel composition $P\parallel Q$ leaves the order of events unspecified. A classical scheduler, even one whose later choices depend on earlier outcomes, yields a process in the classical-control class, hence a causally separable one. The switch is a coherent scheduler, well-typed on every non-signalling pair of channels, whose output applies the two operations in a superposition of orders; it is causally non-separable, and it has content only when the two programs act on a shared target or communicate, since on disjoint systems every order gives the same channel. The analogy is therefore between $\parallel$ and the argument type of a scheduler, not between $\parallel$ and the switch itself; this is what makes \WP{2d} well-posed, and also what makes it hard, since coherent control of order is not an algebraic operation with classical continuations. Finally, Simmons's own outlook \cite{simmons2024thesis} proposes reading send and receive operations as physical processes so as to assign causal types to communicating programs, with causal consistency playing the role of deadlock freedom and the session-type logics BV, GV and MAV as the point of contact; that is, from the physics side, the direction of \WP{2a} and \WP{2b}.

% ---------------------------------------------------------------------
\section{Synthesis: three notions of causality and the gap}\label{sec:synthesis}

Table \ref{tab:synthesis} summarises the review. Each column is one extensional notion of causality; each row asks the same question of it.

\begin{table}[ht]
\centering\small
\begin{tabularx}{\textwidth}{>{\raggedright\arraybackslash}p{3.1cm}XXX}
\toprule
 & \textbf{Intervention} & \textbf{Signalling} & \textbf{Factorisation}\\
\midrule
Intensional datum & dataflow of a program / SCM & temporal order of events: trace $\sigma$, comb order, pomset & DAG\\
Extensional characterisation & contextual equivalence under interventions & signalling types ($\CausC$); causal (non-)separability & conditional independence / d-separation\\
Correspondence theorem & full abstraction for concrete models (\S\ref{sec:intv-wherewestand}) & full abstraction and draft completeness relate three equalities (\S\ref{sec:order-wherewestand}); the signalling reading is Conjecture \ref{conj:b1} & soundness of d-separation, synthetically \cite{fritz2023dseparation}; completeness classically\\
Semantics of programs & first-order and higher-order, acyclic & single channel, classical control & Markov categories; interventions by surgery\\
Next hard case & cyclic SCMs & multiple channels; causally non-separable processes & cyclic and quantum causal models\\
\bottomrule
\end{tabularx}
\caption{Three notions of causality and where the two routes stand.}\label{tab:synthesis}
\end{table}

Three remarks. First, the example of \S\ref{sec:intro} and the asymmetry it exposes: for intervention the correspondence between intensional and extensional structure is a theorem; for signalling there is a semantics with order as first-class data and a conjecture. Second, non-DAG structure is the next hard case in every column, but it is not the same case: feedback in cyclic SCMs, causal non-separability in quantum processes, and unscheduled parallel composition in programs are three different departures from a fixed order, and a unifying account must distinguish them. Third, and this is the gap: there is no framework in which the three columns are structure on one category. Markov categories cover factorisation and, through surgery, part of intervention; $\CausC$ covers signalling for quantum and finite classical processes; quantum causal models relate all three, with a compositional presentation in Lorenz and Tull's string diagrams, but are not formulated as a semantics of programs. The two semantics of \S\ref{sec:intv-wherewestand} and \S\ref{sec:order-wherewestand} cover the first two columns separately, and both are families of morphisms indexed by an intensional datum, intervention grades or traces, with a reindexing action, which is the common structure a generalisation would start from.

Two obstructions should be displayed rather than hidden. The first is \emph{finite dimension}: the $\CausC$ construction needs a compact closed base with invertible dimensions, which $\Mat$ and finite-dimensional $\CP$ provide, but measurable spaces, quasi-Borel spaces and the infinite-biproduct completion $\CPinf$ do not. The bridge of \WP{3a} is therefore stated on the finite fragment, and the comb reading of \WP{2b} is stated per trace, over finite-dimensional $\CP$, with the trace index kept external. The second is \emph{sub-normalisation}: conditioning in $\LHO$ denotes sub-normalised maps, and branching in $\iotq$ makes individual trace components sub-normalised, so neither lands in a first-order $\CausC$ type as it stands; testers and strategies \cite{chiribella2009networks,gutoski2007games} are the objects that accommodate sub-normalised families, Simmons \cite{simmons2024thesis} surveys the ways of admitting partial processes into $\CausC$ and finds each wanting in some respect, and normalisation boxes \cite{jacobsszelesstein2025inference} are the diagrammatic tool for renormalising after conditioning. I also note that the factorisation column is the thinnest for me: it enters through Markov categories and the diagrammatic causal-model literature, not through my own results, and \WP{3a} includes a factorisation reading for that reason.


% =====================================================================
\part{Thesis Proposal}
% =====================================================================

\section*{Notation and background in brief}
\begin{itemize}[nosep]
\item \emph{Interventional languages} $\Lk{1}\subseteq\Lk{2}\subseteq\Lk{3}$ and $\LHO$: first-order probabilistic languages with intervention labels $\ell$ read by $\mathsf{applyIntv}$, interventions $\dok(\ell\mapsto v)\{t\}$ by substitution, twinning by vectorisation, and a higher-order extension. $\equiv_k$: contextual equivalence at rung $k$. Grading: the denotation of an $\Lk{2}$ program is indexed by the interventions available.
\item \emph{Communication theory} $\iotq$: the tensor of an I/O theory ($\opin$ receives a qubit, $\opout$ sends one) with Staton's theory $\mathsf{Q}$ ($\opnew$, $\opmeasure$ with two branches, $\applyU{U}$). Binding: $\applyU{U}(\vec a,\ \vec a.\ t)$ applies $U$ to $\vec a$ and continues as $t$; $\mathsf{discard}(q, t)$ measures and ignores. $\approx_q$: operational equivalence against quantum streams. $T(X)=\bigoplus_\sigma \mathrm{qbit}^{|\sigma|}\otimes X$ on $\CPinf$: the free I/O monad; $\llbracket t\rrbracket_\sigma$: the $\sigma$-component of a denotation.
\item \emph{Signalling}: a \emph{comb} is a process with slots in a fixed order; a \emph{tester} is a sub-normalised family of combs summing to one. $\CausC(\mathcal{C})$: Kissinger and Uijlen's category of higher-order processes over a precausal $\mathcal{C}$, typed by signalling constraints; $\Mat$ and finite-dimensional $\CP$ are (additive) precausal; $\FinStoch$ (stochastic matrices) is the first-order part of $\CausC(\Mat)$; graph types \cite{simmons2024thesis} are the signalling types of an arbitrary DAG. Classical-control and quantum-control circuits: the classes of Wechs et al.\ \cite{wechs2021qcqc}; the switch is in the second and not the first.
\item \emph{Factorisation}: Markov categories, Bayesian networks as string diagrams, d-separation.
\item M1--M3: reading milestones in \S\ref{sec:plan}.
\end{itemize}

\section{Research topic and fundamental challenges}

The thesis will develop a \emph{compositional semantics of causal structure for probabilistic and quantum programs} in which (1) interventional structure is characterised by contextual equivalence, which is done; (2) the order structure of communicating quantum programs is characterised by signalling types, which is the main new technical work; and (3) the two are related by a functorial bridge on their common finite classical fragment. An axiomatic framework in which intervention, signalling and factorisation are all structure on one category is the stretch goal, attempted only once (3) has decided whether it has the right shape. The \emph{minimum viable thesis} is (1), the completeness theorem for $\iotq$, the signalling reading of its single-channel semantics and its extension to independent channels, and the bridge (3).

The fundamental challenges are these. (i) \emph{Finite dimension and sub-normalisation.} Signalling structure is built over compact closed categories of finite-dimensional systems, whereas the interventional semantics is measure-theoretic and the communication semantics uses an infinite-biproduct completion; and both conditioning and measurement branching produce sub-normalised maps, which are not elements of first-order signalling types. (ii) \emph{Parallel composition is not an algebraic operation} \cite{stark:freamp,CSP_van_Glabbeek_2010}, so threads cannot simply be added to $\iotq$ as an operation; and coherent control of order is not an operation with classical continuations at all. (iii) \emph{Completeness and full-abstraction proofs are model-specific and long}: the single-channel completeness proof has taken most of a year and is not yet closed. (iv) \emph{No accepted standard to test against}: the existing behavioural semantics for communicating quantum programs disagree with one another \cite{lqCCS_10.1145/3632885}, so evaluation must be by theorems and by matching known protocols and known physics.

\section{Goal}

I want a mathematical account of what it means for one part of a program to cause another, that works the same way whether ``cause'' means ``if I change this, that changes'', ``this can send a message to that'', or ``these two are independent given those'', and whether the program is classical-probabilistic or quantum.

\section{State of the art and why the problem is hard}

The existing notions of equivalence for causal models (Bongers et al.; causal spaces; causal abstraction) say when two models have the same interventional behaviour, but they are not equivalences of \emph{programs}: there is no notion of context, and no higher order. Causal probabilistic programming languages (OmegaC, ChiRho, MultiVerse) implement interventions and counterfactuals at higher order, but by evaluation strategies, without a denotational or equational account of what they compute. Quantum process calculi (CQP, qCCS, lqCCS) are the only semantic treatment of communicating quantum programs; they are operational, their bisimilarities disagree, and they do not separate communication from concurrency. The theory of higher-order quantum processes (combs, process matrices, the $\CausC$ construction) has a precise account of signalling and of definite versus indefinite order, but it has never been connected to a programming language. Markov categories, the synthetic home of factorisation, have no primitive notion of intervention or of signalling.

The problem is hard for three reasons. Each notion has been developed with its own mathematics, and the natural meeting points---parallel composition, cycles, indefinite order---are precisely where the standard tools (algebraic theories, acyclic graphs, fixed-order combs) stop working. The bases do not match: signalling needs finite dimension and compact closure, the interventional semantics is measure-theoretic, and conditioning and branching produce sub-normalised maps. And the correspondence theorems that would relate the notions are each a substantial proof.

\section{Approach}\label{sec:approach}

\paragraph{Method.} The method is the one that has now worked twice. (1) Fix a small language whose \emph{syntax} carries the intensional causal structure: labels for interventions, I/O operations for communication events. (2) Define an \emph{extensional} equivalence from what an observer may do: quantify over interventional contexts, or over quantum-stream environments. (3) Build a \emph{denotational} model in which the intensional structure becomes an index, a grade or a trace, with a reindexing action. (4) Prove a \emph{correspondence theorem}, full abstraction or completeness, between (2) and (3). Work packages 1 and 2 push each existing instance to its next hard case; work package 3 asks for the common theory.

\paragraph{What is new.} Three things. Causal structure is \emph{defined} by a contextual equivalence of programs, rather than by a property of a graph; this is what makes it compositional and higher-order. The order of communication events is \emph{denotational data}, an index of the model, rather than a label on an operational transition; this is what makes a signalling reading possible. And the bridge between the two is a functor between a graded category and a signalling-typed category, which, if it exists, says that intervening at a hole and plugging a channel into a slot are the same operation seen from two sides.

\paragraph{Two worked examples.} The two examples below are the same statement in two guises, and they are the first thing each work package must preserve.

\emph{Example A (interventional).} In $\Lk{2}$, interventions are exposed as typed free variables $\ell$, the intervention labels, read through $\mathsf{applyIntv}$, which returns the label's value when it is set and the program's own value otherwise. Consider a mediation model:
\begin{lstlisting}
let t  = bernoulli 0.5 in              -- treatment
let t' = applyIntv $\ell$ t in             -- intervenable copy of t
let m  = bernoulli (f t') in           -- mediator
let y  = bernoulli (g t' m) in         -- outcome
let z  = bernoulli 0.3 in              -- unrelated variable
return (y, z)
\end{lstlisting}
An intervention is substitution for the label: under $\dok(\ell\mapsto\mathsf{Just}\ \mathsf{true})\{-\}$ the variable $t'$ denotes the constant $\mathsf{true}$, and the total effect of $t$ on $y$ is the difference between the distributions of $y$ under $\dok(\ell\mapsto\mathsf{Just}\ \mathsf{true})$ and $\dok(\ell\mapsto\mathsf{Just}\ \mathsf{false})$. The natural direct effect---the effect of $t$ on $y$ with the mediator held at the value it \emph{would have taken} under $t=\mathsf{false}$---is a cross-world quantity; in $\LHO$ it is expressed by twinning the program and intervening on the mediator of one world with a value read from the other. The invariance claim is that $z$ has the same distribution in every interventional context: $z$ is causally independent of $t$ because the dataflow does not connect them, and $\equiv_2$ turns this into an equation between programs of the same type, namely that the program above with $\mathtt{return}\ z$ in place of $\mathtt{return}\ (y,z)$ is $\equiv_2$ to $\mathtt{let}\ t' = \mathtt{applyIntv}\ \ell\ \mathtt{true}\ \mathtt{in}\ \mathtt{let}\ z = \mathtt{bernoulli}\ 0.3\ \mathtt{in}\ \mathtt{return}\ z$, in which the label is consumed but ignored. Full abstraction says that the graded denotation detects exactly these equations.

\emph{Example B (communication).} In $\iotq$, $\opin: \langle 0 \mid 1\rangle$ receives a qubit and binds it, $\opout: \langle 1\mid 0\rangle$ sends one, and the operations of $\mathsf{Q}$ are $\opnew$, $\opmeasure$ (which branches on the outcome) and $\applyU{U}$. Consider the term that receives a qubit, prepares a Bell pair locally, performs a Bell measurement on the received qubit and one half of the pair, applies the Pauli correction determined by the two outcomes to the other half, and sends it:
\begin{align*}
\mathit{tele} \;=\; \opin\big(q.\ \opnew(a.\ \opnew(b.\ &\applyU{H}(a,\ a.\ \applyU{CX}((a,b),\ (a,b).\\
&\applyU{CX}((q,a),\ (q,a).\ \applyU{H}(q,\ q.\ \opmeasure(q,\ c_{0},\ c_{1})))))))\big)
\end{align*}
where $c_{i} = \opmeasure(a,\ \applyU{Z^{i}}(b,\ b.\ \opout(b, x)),\ \applyU{Z^{i}X}(b,\ b.\ \opout(b, x)))$. Write $T[y]$ for the local part of this term with the continuation $y(b)$ in place of each $\opout(b,x)$; $T[y]$ is a term of $\mathsf{Q}$ alone, its denotation is the identity channel from $q$ to $b$, and Staton's theory is complete \cite{statonpopl15}, so $T[y] = y(q)$ is derivable. Equational derivations are closed under substitution for continuation variables, hence $\mathit{tele} = \opin(q.\ \opout(q, x))$ in $\iotq$: the term is provably equal to a bare relay. Its denotation has a single non-zero trace component, at $\sigma=(\opin,\opout)$, and that component is, up to the Choi correspondence, the identity on one qubit: a comb with one input slot followed by one output slot.

The derivation above uses only the local theory; the \emph{tensor} is exhibited by the following equation from our paper, which holds in $\iotq$ but not in the sum of the two theories:
\begin{align*}
&\opnew(q.\ \applyU{H}(q,\ q.\ \opnew(p.\ \applyU{CX}((q,p),\ (q,p).\ \opout(p,\ \mathsf{discard}(q, x))))))\\
&\quad = \opnew(q.\ \applyU{H}(q,\ q.\ \opmeasure(q,\ \opnew(p.\ \opout(p, x)),\ \opnew(p.\ \applyU{X}(p,\ p.\ \opout(p, x)))))).
\end{align*}
Preparing a Bell pair, sending one half and discarding the other is indistinguishable, to any environment, from tossing a fair coin and sending its outcome as a classical bit. The derivation needs the commutation axiom $\opmeasure(a,\ \opout(b,x),\ \opout(b,y)) = \opout(b,\ \opmeasure(a,x,y))$: the measurement of the discarded half and the sending of the other half may be reordered because neither can signal to the other. This is the communication-side counterpart of the invariance claim in Example A.

The teleportation term has one branch and one trace, which hides what makes the signalling reading non-trivial. A term that branches on a measurement, such as $\opnew(a.\ \opout(a,\ \opin(q.\ \opmeasure(q,\ x,\ \opnew(b.\ \opout(b, x))))))$, has an $(\opout,\opin)$-component that is post-selected on the later measurement outcome: it is sub-normalised, and on its own it is not a comb, since its trace depends on the later input. Only the family of components over all traces, with the later wires discarded, satisfies the comb conditions. The right objects are therefore testers and strategies \cite{chiribella2009networks,gutoski2007games}, and the conjecture is:
\begin{conj}\label{conj:b1}
For every closed term $t$ of $\iotq$ and every trace $\sigma$, the component $\llbracket t\rrbracket_\sigma$ is a probabilistic comb with slot order $\sigma$, and the family $(\llbracket t\rrbracket_\sigma)_\sigma$ is a variable-length tester: for each $\sigma$, the sum of the components over all extensions of $\sigma$, with the later wires discarded, is a deterministic comb with slot order $\sigma$.
\end{conj}
\noindent Stated this way the conjecture is expected to follow by structural induction on terms, since sequencing is comb composition, and its content is the normalisation bookkeeping across traces of different lengths; the question that is genuinely open is the characterisation in Conjecture \ref{conj:b2}.

\paragraph{Work package 1: interventional route---equational theory and cycles.}
\begin{itemize}[nosep]
\item[1a] \textit{(medium risk)} Relate the equational theory of contextual equivalence to the normalisation-box equational theories for causal Bayesian networks \cite{jacobsszelesstein2025inference,agrawal2026equational}, and settle whether $\equiv_3$ is strictly finer than two-world counterfactual equivalence. \emph{Done:} a coincidence theorem on the programs definable from causal Bayesian networks, or a separating example; and the strictness answer. \emph{Why it may succeed:} both are equational theories over the same string-diagram semantics. This item is in the plan because it is cheap and because its outcome feeds \WP{3c}.
\item[1b] \textit{(high risk; medium if restricted to finite cyclic models with unique solutions)} Cyclic SCMs and feedback: identify the fixpoint construct whose denotation is Bongers et al.'s solution function, and what the graded semantics must become to support it. \emph{Done:} a language extension with a sound semantics for the restricted class, and Example A with a feedback edge. \emph{Why it may succeed:} the solution-function semantics is already a fixpoint; the question is which language construct denotes it. Outside the minimum viable thesis.
\end{itemize}

\paragraph{Work package 2: communication route---multiple channels and signalling.}
\begin{itemize}[nosep]
\item[2a] \textit{(medium risk; in the minimum viable thesis)} Independent channels in a sequential program: $\iotq$ with operations $\opin_c$, $\opout_c$ for each channel $c$, where events on different channels commute, so that the monad is indexed by Mazurkiewicz traces over a partially commutative monoid. \emph{Done:} the theory defined, non-degeneracy by a multi-stream operational model, soundness and adequacy of the trace-indexed denotational model; completeness by re-running the single-channel proof, whose reorganisation step and factorisation lemma are channel-agnostic. \emph{Why it may succeed:} the proof architecture is unchanged; what is new is the commutation of independent I/O events, which the quantum-stream model must validate. \emph{Test:} teleportation as a two-party program; entanglement swapping.
\item[2a$'$] \textit{(high risk; stretch)} Threads: tensor $\iotq$ with the dynamic-threads theory of Kammar et al.\ \cite{kammar2026threads} and ask which of its equations survive the quantum setting, in the manner of the Eckmann--Hilton analysis in my dissertation; traces become pomsets. \emph{Done:} the theory defined and shown non-degenerate; a soundness theorem.
\item[2b] \textit{(low risk for Conjecture \ref{conj:b1}; medium for Conjecture \ref{conj:b2})} Which fragment of higher-order quantum computation does $\iotq$ axiomatise?
\begin{conj}\label{conj:b2}
For each trace $\sigma$, the component $\llbracket t\rrbracket_\sigma$ lands in the comb type for $\sigma$ in $\CausC(\CP)$, the branch family is typed by the coproduct and probabilistic-choice operators of the additive construction \cite{simmons2024thesis}, and the assignment is functorial in sequencing. Every deterministic comb of finite depth with qubit slots is the component of some term, and the equations of $\iotq$ are exactly those validated by the comb types.
\end{conj}
\noindent The first sentence is the typing statement; the second is expected to hold by the comb realisation theorem \cite{chiribella2009networks} together with the universality of $\mathsf{Q}$; the third, completeness of $\iotq$ relative to $\CausC$, is the open question, and causal logic \cite{simmons2024thesis,simmons2024completelogic} gives it a sharper form: which fragment of causal logic do the I/O orders of $\iotq$ terms realise, and is every causally consistent wiring in that fragment denotable? The trace index is kept external, so the construction is over finite-dimensional $\CP$, which is additive precausal, and not over $\CPinf$, which fails the finite-causal-basis axiom. \emph{Done:} \ref{conj:b1} proved; the typing functor constructed; the realisation statement proved or refuted; relative completeness answered or reduced to a precisely stated question. \emph{Why it may succeed:} the objects on both sides are finite-dimensional CP maps with the same wire structure, and the single-channel completeness proof already produces the normal forms that a comparison with comb types needs.
\item[2c] \textit{(high risk; conditional on 2a and 2b)} Distributed quantum computing through process matrices. In the process-matrix picture the local programs are combs and the network that wires them is a process matrix; the conjecture is that the classical-control and quantum-control classes of Wechs et al.\ \cite{wechs2021qcqc} are the right denotational universe for the \emph{network} component of a distributed program, with local programs denoting as in \WP{2b}, so that ``realisable as a circuit with controlled order'' becomes a decidable predicate on denotations; it cannot be a $\CausC$ type, since every process matrix is an affine combination of definite-order ones \cite{simmons2024thesis}. Two hypotheses, in increasing order of content:
\begin{hyp}\label{hyp:sched}
A parallel composition of $\iotq$ programs under a classical scheduler whose choices may depend on earlier outcomes denotes a process in the classical-control class, hence a causally separable one; an outcome-independent scheduler denotes a convex mixture of fixed-order combs.
\end{hyp}
\begin{hyp}\label{hyp:conv}
Conversely, every process in the classical-control class is the denotation of some parallel composition under some classical scheduler.
\end{hyp}
\noindent Hypothesis \ref{hyp:sched} is close to the definition of the classical-control class and is a formalisation exercise; Hypothesis \ref{hyp:conv} would be a result. \emph{Done:} the translation from parallel compositions to process matrices with Hypothesis \ref{hyp:sched} proved; Hypothesis \ref{hyp:conv} proved or refuted; and, for each protocol in a catalogue (teleportation, entanglement swapping, a BB84 fragment, a two-party comb), a term whose denotation is the known process operator. The switch and the process of Oreshkov, Costa and Brukner are not in the classical-control class, so Hypothesis \ref{hyp:sched} already shows that no such composition denotes them; they are the boundary of the model, not independent tests. \fillin{the Kim reference for this item} \emph{Why it may succeed:} the physics side has already characterised the target classes, and the translation reuses the combs of \WP{2b}.
\item[2d] \textit{(high risk, exploratory; after 2c)} Generalised switches as quantum schedulers. A pomset over channel events has a set of linear extensions, and a coherent superposition over a subset of the permutations of $N$ single-use nodes is the $N$-party switch with a restricted control state, which lies in the quantum-control class; realisability of pomset-indexed switches is therefore expected to be a corollary rather than a milestone. The questions with content are: (a) which quantum-control processes arise as schedulers over linear extensions, since that class also allows the order to be updated coherently by the operations' outputs, which is strictly more; (b) what a ``quantum scheduler'' primitive in a programming language is, given that coherent control of order is not an algebraic operation with classical continuations, and how it relates to the PBS-calculus \cite{clement2020pbs}, which has the switch as a term; (c) the well-posedness condition that the composed programs act on a shared target or communicate, with identical wiring in every branch so that the control does not leak. \emph{Done:} the definition; the characterisation (a) or a precise separation; a proposal for (b) with a sound semantics for the two-party case recovering the switch.
\end{itemize}

\paragraph{Work package 3: a unifying categorical account.}
\begin{itemize}[nosep]
\item[3a] \textit{(medium risk; in the minimum viable thesis)} A bridge theorem on the common finite classical fragment. Finite stochastic processes live in both worlds: as morphisms of $\FinStoch$, over which the interventional semantics can be instantiated, and as the first-order part of $\CausC(\Mat)$. The reframing that gives the theorem content is that an $\Lk{2}$ program with intervention labels \emph{is} a comb whose slots are the labels: a perfect intervention plugs a state into a slot, a soft intervention plugs a channel, and the tightest $\CausC$ type of a program is determined by its dataflow (tensor for labels with no dataflow between them, sequential for ordered ones). The theorem to prove is that this type assignment is a functor from the graded interventional category into $\CausC(\Mat)$, compatible with composition and with the intervention grades, and that it reflects interventional equivalence; for a single label the two sides coincide by unfolding definitions, and the content is in composition and in the higher-order fragment. Two refinements are built in. First, perfect and soft interventions differ: invariance of $Y$ under all perfect interventions at a slot does not imply invariance under all soft ones, whereas invariance under all soft interventions is exactly no-signalling from the slot; the theorem is stated for soft interventions, and the perfect case is reported as the weaker relation it is. Second, a \emph{factorisation} reading is included: the conditional independences of the Markov-category image, via Fritz and Klingler's synthetic d-separation \cite{fritz2023dseparation}, as the third face of the same functor. A further hypothesis, to be checked early because it would make the interventional route literally a fragment of the communication route: an intervention label is an $\opout$/$\opin$ pair with the intervener as the environment, $\dok(\ell\mapsto v)$ being the stream that ignores the output and returns $v$ and ``no intervention'' the echo stream, so that $\equiv_2$ is equivalence under all memoryless echo-or-replace streams. The higher-order fragment is obstructed by sub-normalisation from conditioning, so the functor is first defined on the normalised fragment and the extension to testers is a stated sub-question. \emph{Done:} the functor and its properties, or a counterexample program; the soft/perfect separation; the factorisation reading. Either outcome is a result: a counterexample is a publishable separation between interventional and signalling independence. Comb disintegration \cite{jacobs2019causal} is the prior work this item must be positioned against; it uses combs to compute with a fixed classical model, whereas here the comb type is assigned to a program and compared with its contextual equivalence.
\item[3b] \textit{(risk unknown until reading milestone M2)} Recast quantum causal models \cite{barrett2019quantumcausalmodels} in the graded setting over finite-dimensional $\CP$, with nodes as labels. The process operator of a quantum causal model is not in general a comb, since nodes with no path between them form a non-signalling rather than a sequential shape, so the graded category must accommodate DAG-shaped signalling types, for which Simmons's graph types \cite{simmons2024thesis} are the vocabulary. \emph{Done:} the quantum Markov condition recovered as a property of the intervention-graded category, or a precise statement of what is missing. \emph{Why it may succeed:} the quantum causal model's interventions are local instruments at nodes, which is what intervention labels denote.
\item[3c] \textit{(high risk, stretch; only after 3a)} An axiomatic framework, in the spirit of Markov categories, in which intervention, signalling and factorisation are structures with stated relationships, instantiated by $\FinStoch$ and finite-dimensional $\CP$, such that the existing correspondence theorems become instances of general ones. \emph{Done:} axioms, the two instances, and at least one general theorem specialising to a known one. \emph{Why it may succeed:} both existing semantics are families indexed by an intensional datum with a reindexing action, so the generalisation starts from two worked instances rather than from an abstraction exercise, and 3a decides within one term whether it has the right shape.
\end{itemize}

\section{Evaluation}

The thesis is theoretical, and its results are evaluated as theorems with the ``done'' criteria stated above: completeness, full abstraction, the typing and realisation statements of \WP{2b}, and the coincidence or separation results of \WP{3a}. Case studies have pass/fail criteria: for the protocol catalogue of \WP{2c}, each protocol either has a term whose denotation is the known process operator or a proof that none does; for the causal-inference examples (mediation, instrumental variables, feedback), each is reproduced in the executable semantics and gives the textbook answer. The Haskell implementation of $\LHO$ \cite{ARTIFACT_wang_2026_23207221} will be extended to each new construct, so that every semantic definition remains executable.

\section{Impact and risks}

If successful, the thesis gives denotational foundations for causal probabilistic programming, and hence for causal machine-learning languages, in which interventions are mathematical operations rather than evaluation strategies; a semantic model of distributed quantum programming in which realisability as a controlled-order circuit is a predicate on denotations; and a dictionary between causal inference and quantum causality, two communities that use the same word for different things.

\begin{table}[ht]
\centering\small
\begin{tabularx}{\textwidth}{lcX}
\toprule
Item & Risk & Fallback\\
\midrule
1a equational theories & medium & a separating example is itself a result\\
1b cyclic models & high & restrict to finite, unique-solution models; outside the minimum viable thesis\\
2a independent channels & medium & soundness and adequacy without completeness\\
2a$'$ threads & high, stretch & dropped without loss\\
2b comb typing (B1) & low & ---\\
2b relative completeness (B2) & medium & the typing functor and realisation alone characterise the fragment\\
2c network semantics & high, conditional & Hypothesis \ref{hyp:sched} and the catalogue alone\\
2d generalised switches & high, conditional & dropped without loss\\
3a classical bridge & medium & a counterexample is a publishable separation\\
3b quantum causal models & unknown until M2 & a precise statement of what is missing\\
3c axiomatic framework & high, stretch & 3a becomes the unifying chapter\\
\bottomrule
\end{tabularx}
\caption{Risks and fallbacks.}
\end{table}

The generic risk is breadth. The mitigation is the \emph{minimum viable thesis}: the interventional chapter (done), the communication chapter (the completeness theorem, with the MSc results as background), the signalling reading of \WP{2b} and the independent-channel extension of \WP{2a}, and the bridge theorem of \WP{3a}. Every item in it is rated low or medium. The remaining items are attempted only when their preconditions are met.

\section{Plan of work, milestones, and training}\label{sec:plan}

I am in my fourth term. The baseline schedule is the minimum viable thesis; conditional items are marked.

\begin{table}[ht]
\centering\small
\begin{tabularx}{\textwidth}{lX}
\toprule
Term & Milestones\\
\midrule
MT26 (4) & Transfer of status. POPL revision. Reading milestone M1: Simmons's thesis \cite{simmons2024thesis}, chapters 2--3, and the papers of \S\ref{sec:concurrency}--\ref{sec:hoqc}.\\
HT27 (5) & Completeness manuscript closed and submitted. Conjecture \ref{conj:b1} proved. Reading milestone M2: quantum causal models \cite{barrett2019quantumcausalmodels}.\\
TT27 (6) & \WP{2b}: typing functor and realisation. \WP{3a}. \emph{Deliverable:} paper on completeness and the comb reading; 3a answered.\\
MT27--HT28 (7--8) & \WP{2a}: independent channels. \WP{1a}. Reading milestone M3: concurrency semantics in depth, before any work on 2a$'$. \emph{Deliverable:} paper on the classical bridge or on 2a. \textbf{Confirmation of Status in HT28}, with the minimum viable thesis in hand.\\
TT28 (9) & Buffer for Confirmation of Status. \WP{3b}. Start of \WP{3c} if 3a succeeded.\\
MT28 (10) & \WP{2c} if 2a and 2b are complete, otherwise \WP{3c}; \WP{1b} on the restricted class if time permits.\\
HT29--TT29 (11--12) & \WP{2d} only if 2c is done by MT28; otherwise consolidation. Thesis writing; submission by the end of TT29.\\
\bottomrule
\end{tabularx}
\end{table}

Consolidation items, not research: an implementation of $\LHO$ beyond MonadBayes if time permits, and a denotational account of OmegaC's core as a worked example. Dissemination targets are POPL, LICS, FoSSaCS, MFPS and QPL.

The thesis will be organised as: 1 Introduction; 2 Background (Part I of this document); 3 Interventional semantics of causal probabilistic programs \fillin{delineate your contribution from coauthors'}; 4 Order semantics of quantum communication, consisting of the completeness theorem and \WP{2a}, with the MSc dissertation as background and not as DPhil material \fillin{confirm reuse rules with supervisor}; 5 Signalling structure and higher-order processes (\WP{2b}--2d); 6 The bridge and the unifying account (\WP{3}); 7 Conclusion.

\paragraph{Training and engagement.} \fillin{required by the assessors' criteria: seminars attended (PL group, quantum group); courses audited; the PLanQC 2025 talk, with a paragraph on it; teaching; planned activities for the coming year}.

\section{Statement of the subject of the thesis}

Compositional semantics of causal structure in probabilistic and quantum programs: interventional, signalling and factorisation characterisations, and the correspondences between them.


\printbibliography

\end{document}


% =====================================================================
% SCRUTINY LOG (draft.tex)
% =====================================================================
% Round 1 (self-review against the source files, 3 Nov 2026):
%  - The lqCCS quotation supplied by an earlier examiner pass is not in
%    that paper; replaced by a paraphrase of its abstract.
%  - The intro's commutation axiom was schematic; replaced by axiom (V)
%    of the FoSSaCS paper, with (Z) used in Example B.
%  - Example B credited the tensor axioms for a derivation that uses only
%    Q-completeness plus substitution; the tensor is now exhibited by the
%    paper's own Bell-pair/coin equation (an instance of (Z)).
%  - Example A rewritten in the paper's label/applyIntv/do syntax.
%  - Confirmed: L_HO widens to the unnormalised measure monad for
%    conditioning; the first-order semantics is over Prob.
% Round 2 (two independent examiner reads, 3-5 Nov 2026):
%  PL-semantics examiner, verified against sources and applied:
%   twinning is L3 -> L2, not into L3; the Bongers correspondence is on
%   compiled L3 programs (obs/intv/2-world); full abstraction needs a
%   concrete model; L3 -> L_HO is semantics-preserving, not conservative;
%   the intro's <=_2 equation was ill-typed (labels are linear) and is now
%   between programs of the same type; the MSc results are not first-year
%   work (abstract and S1 fixed); the operational equivalence is not
%   contextual (said, with the open question); completeness moved to term
%   5; milestones renumbered; Confirmation of Status in term 8 with 9 as
%   buffer; 2a split into independent channels (Mazurkiewicz traces,
%   medium, in the MVT) and threads (high, stretch); risk table and MVT
%   now agree; "what is new" paragraph and a notation box added to Part
%   II; S9 made self-contained; the factorisation face given a deliverable
%   in 3a; the intervener-as-environment hypothesis added to 3a; comb
%   disintegration, centrality, Mazurkiewicz traces added to the review;
%   joint-work delineation and training left as FILLs.
%  Quantum-foundations / causal-inference examiner, applied:
%   Conjecture B1 was false as stated (branch components are
%   post-selected); restated as a tester/strategy property with the
%   counterexample shown; B2 restated over finite-dimensional CP with the
%   trace external, realisation and relative completeness separated; the
%   "unnormalised => not precausal" obstruction was wrong (Mat(R>=0) is
%   Kissinger-Uijlen's own precausal example) and is replaced by finite
%   dimension / compact closure and sub-normalisation; Wechs et al.
%   stated exactly (three nested classes; switch in quantum-control but
%   not classical-control; classical-control => separable; quantum-control
%   => causal correlations; exhaustiveness is a proposal); causal
%   inequalities and the switch/OCB distinction added; realisability of
%   the switch hedged and cited; Caus types corrected (the indefinite-order
%   type contains all process matrices; separability is not a type);
%   Simmons-Kissinger: "additives" dropped; quantum causal models
%   described at the level of their definitions; cyclic models hedged;
%   CDP 2010 cited for the causality axiom and Coecke's result stated
%   correctly; Fritz-Klingler: soundness only; conditionals as an extra
%   axiom; abduction-action-prediction; SWIGs single-world caveat;
%   Lorenz-Tull attribution fixed; 2c reframed (programs are combs, the
%   network is the process matrix; Hyp. 1 near-definitional, Hyp. 2 the
%   result; "realisable" is a predicate, not a type; switch and OCB are
%   the boundary); 2d reframed (realisability a corollary of the N-switch;
%   "exactly quantum-control" dropped; coherent control not algebraic;
%   PBS-calculus cited; shared-target condition); 3a reframed (labels as
%   comb slots; soft vs perfect; FinStoch).
% Addition (5 Nov 2026): Simmons's DPhil thesis (Causal types for higher-order
%  quantum theory, Oxford 2024) worked into S4.4, S4.5, S5, the notation box,
%  B2, 2c, 3b and the reading plan; it supplies the additives, graph types,
%  partiality survey, the affine-universality corollary and causal logic.
%  Not applied: the suggestion to drop "factorisation" from the title
%   (kept, with a deliverable in 3a instead); experimental-switch papers
%   beyond one citation.
